\section{Forest Variants}

\begin{frame}{From Random Forests to Other Forest Variants}
\begin{itemize}
  \item Random forests average many randomized trees to reduce variance.
  \pause
  \item The same recipe -- random partitioning, repeated many times -- is useful beyond classification and regression.
  \pause
  \item Today: three variants that reuse the forest idea for new purposes.
\end{itemize}

\pause
\begin{rlintuition}
A forest is not just a predictor. It is a way to summarize how data is partitioned by randomized splits, and that summary itself is useful information.
\end{rlintuition}
\end{frame}

\begin{frame}{Three Ways to Reuse Random Partitioning}
\begin{itemize}
  \item \textbf{Isolation Forest} -- use how \emph{quickly} a point gets isolated as an anomaly score.
  \pause
  \item \textbf{Extremely Randomized Trees (Extra-Trees)} -- push randomization further into the split itself, not just the data.
  \pause
  \item \textbf{Quantile Regression Forests} -- keep the full leaf distribution instead of only its mean.
  \pause
  \item All three reuse bagging and/or feature randomness from Lecture 6.
\end{itemize}
\end{frame}

\begin{frame}{Isolation Forest: Anomalies Are Easy to Isolate}
\begin{columns}[T]
\column{0.53\textwidth}
\begin{itemize}
  \item Idea: anomalies are few and different.
  \pause
  \item Random axis-aligned cuts isolate them in a small box quickly.
  \pause
  \item Normal points sit in dense regions and need many cuts to separate from their neighbors.
  \pause
  \item So: \emph{short average path to a leaf} $\Rightarrow$ likely anomaly.
\end{itemize}

\pause
\column{0.47\textwidth}
\begin{rlintuition}
Isolation Forest never models what ``normal'' looks like. It only measures how easy a point is to cut off from the rest.
\end{rlintuition}
\end{columns}
\end{frame}

\begin{frame}{Building an Isolation Tree}
\begin{itemize}
  \item Draw a random subsample of size $\psi$ from the data (no class labels needed).
  \pause
  \item At each node, pick a feature $j$ uniformly at random.
  \pause
  \item Pick a split threshold $t$ uniformly at random between the min and max of feature $j$ at that node.
  \pause
  \item Split the node into two children and recurse.
  \pause
  \item Stop when a node has one point left, or a height limit is reached.
  \pause
  \item Repeat for $B$ trees, each on a fresh random subsample, to form the isolation forest.
\end{itemize}
\end{frame}

\begin{frame}{Path Length as an Anomaly Signal}
\centering
\includegraphics[width=0.75\textwidth]{figures/forest_variants/isolation_partition.pdf}

\vspace{0.35em}
\begin{itemize}
  \item Left: the anomaly sits alone, so two random cuts isolate it.
  \pause
  \item Right: the normal point is surrounded by neighbors, so isolating it takes more cuts.
\end{itemize}
\end{frame}

\begin{frame}{Path Length in the Tree}
\centering
\includegraphics[width=0.9\textwidth]{figures/forest_variants/path_length_tree.pdf}

\vspace{0.35em}
\begin{itemize}
  \item The path length $h(x)$ is the number of edges from the root to the leaf containing $x$.
  \pause
  \item Anomalies have small $h(x)$; normal points have large $h(x)$.
\end{itemize}
\end{frame}

\begin{frame}{Worked Numerical Example: Path Lengths}
\small
\begin{itemize}
  \item Subsample size $\psi=16$ for each isolation tree.
  \item Across $B=5$ trees, point $A$ (the anomaly) has path lengths
  \[
    h^{(1)}=2,\quad h^{(2)}=3,\quad h^{(3)}=2,\quad h^{(4)}=1,\quad h^{(5)}=2,
    \qquad \mathbb{E}[h(A)] = 2.0.
  \]
  \pause
  \item Point $B$ (a normal point) has path lengths
  \[
    h^{(1)}=8,\quad h^{(2)}=5,\quad h^{(3)}=6,\quad h^{(4)}=6,\quad h^{(5)}=5,
    \qquad \mathbb{E}[h(B)] = 6.0.
  \]
  \pause
  \item $A$ is isolated three times faster than $B$ on average.
\end{itemize}
\end{frame}

\begin{frame}{Anomaly Score Formula}
\begin{itemize}
  \item Average path length over an unsuccessful binary search tree of $\psi$ nodes:
  \begin{equation}
    c(\psi) = 2H(\psi-1) - \frac{2(\psi-1)}{\psi},
    \qquad H(i) \approx \ln(i) + 0.5772,
    \label{eq:avg-path-length}
  \end{equation}
  where $H(i)$ is the $i$-th harmonic number.
  \pause
  \item The anomaly score for a point $x$ is
  \begin{equation}
    s(x,\psi) = 2^{-\,\mathbb{E}[h(x)]/c(\psi)}.
    \label{eq:anomaly-score}
  \end{equation}
  \pause
  \item $s \to 1$: short paths, likely anomaly.
  \item $s \to 0.5$ or below: long paths, likely normal.
\end{itemize}
\end{frame}

\begin{frame}{Worked Numerical Example: Anomaly Scores}
\small
\begin{itemize}
  \item For $\psi=16$: $H(15) \approx 3.3182$, so by Eq.~\eqref{eq:avg-path-length}
  \[
    c(16) = 2(3.3182) - \frac{2(15)}{16} \approx 4.7614.
  \]
  \pause
  \item Point $A$ (anomaly), $\mathbb{E}[h(A)] = 2.0$:
  \[
    s(A,16) = 2^{-2.0/4.7614} = 2^{-0.420} \approx 0.747.
  \]
  \pause
  \item Point $B$ (normal), $\mathbb{E}[h(B)] = 6.0$:
  \[
    s(B,16) = 2^{-6.0/4.7614} = 2^{-1.260} \approx 0.418.
  \]
  \pause
  \item $A$ scores well above $0.5$; $B$ scores well below $0.5$.
\end{itemize}
\end{frame}

\begin{frame}{Reading the Anomaly Score}
\centering
\includegraphics[width=0.68\textwidth]{figures/forest_variants/anomaly_score_curve.pdf}

\vspace{0.35em}
\begin{itemize}
  \item Score decreases smoothly as the (normalized) average path length grows.
  \pause
  \item A common rule: flag $x$ as an anomaly if $s(x,\psi)$ exceeds a chosen threshold (e.g.\ $0.6$--$0.7$).
\end{itemize}
\end{frame}

\begin{frame}{Isolation Forest: When It Works}
\begin{columns}[T]
\column{0.5\textwidth}
\textbf{Works well when}
\begin{itemize}
  \item anomalies are rare and few,
  \pause
  \item anomalies are ``different'' along some feature,
  \pause
  \item no labeled anomalies are available,
  \pause
  \item linear-time scoring is needed.
\end{itemize}

\pause
\column{0.5\textwidth}
\textbf{Can fail when}
\begin{itemize}
  \item anomalies form their own dense cluster,
  \pause
  \item normal data has many irrelevant features,
  \pause
  \item anomaly notion requires labels.
\end{itemize}
\end{columns}
\end{frame}

\begin{frame}{Extremely Randomized Trees: Pushing Randomization Further}
\begin{itemize}
  \item Random forests already randomize: bootstrap samples + feature subsampling.
  \pause
  \item But the split threshold $t$ within the chosen feature is still found by an exhaustive search for the best impurity decrease.
  \pause
  \item Extra-Trees randomizes the threshold too.
  \pause
  \item Typically trains on the full dataset (no bootstrap), relying on threshold randomness for decorrelation instead.
\end{itemize}
\end{frame}

\begin{frame}{Random Thresholds Instead of Optimal Splits}
\centering
\includegraphics[width=0.95\textwidth]{figures/forest_variants/extra_trees_vs_rf.pdf}

\vspace{0.35em}
\begin{itemize}
  \item Random forest: scan many candidate thresholds, keep the one with the largest $\Delta I$.
  \pause
  \item Extra-Trees: for each candidate feature, draw \emph{one} threshold uniformly at random; pick the best feature among those random draws.
\end{itemize}
\end{frame}

\begin{frame}{Extra-Trees vs Random Forest}
\begin{columns}[T]
\column{0.5\textwidth}
\textbf{Random Forest}
\begin{itemize}
  \item Bootstrap sampling.
  \pause
  \item Exhaustive threshold search per feature.
  \pause
  \item Slower to train.
  \pause
  \item Slightly lower bias.
\end{itemize}

\pause
\column{0.5\textwidth}
\textbf{Extra-Trees}
\begin{itemize}
  \item Usually full dataset, no bootstrap.
  \pause
  \item One random threshold per feature.
  \pause
  \item Faster to train.
  \pause
  \item Slightly higher bias, often lower variance.
\end{itemize}
\end{columns}
\end{frame}

\begin{frame}{Why Does Extra Randomization Help?}
\begin{itemize}
  \item Random thresholds make individual trees weaker (more biased).
  \pause
  \item But they also make trees less correlated with each other.
  \pause
  \item Averaging many weak, decorrelated trees can still reduce variance more than averaging fewer, stronger, correlated trees.
  \pause
  \item Skipping the threshold search also makes training noticeably faster.
\end{itemize}
\end{frame}

\begin{frame}{Quantile Regression Forests: Beyond Point Predictions}
\begin{columns}[T]
\column{0.53\textwidth}
\begin{itemize}
  \item A standard regression forest predicts only the conditional mean $\mathbb{E}[y\mid x]$.
  \pause
  \item It discards information: every leaf actually stores several training $y$-values, not just their average.
  \pause
  \item Quantile Regression Forests keep that full leaf distribution.
\end{itemize}

\pause
\column{0.47\textwidth}
\begin{rlintuition}
A mean prediction tells you the center. A quantile forest tells you the spread -- which is exactly what you need for a prediction interval.
\end{rlintuition}
\end{columns}
\end{frame}

\begin{frame}{Leaf Distributions Instead of Leaf Means}
\centering
\includegraphics[width=0.85\textwidth]{figures/forest_variants/quantile_forest_leaf.pdf}

\vspace{0.35em}
\begin{itemize}
  \item Each leaf holds the $y$-values of the training points that land there.
  \pause
  \item Pooled across all trees, these values approximate the conditional distribution of $y$ given $x$.
\end{itemize}
\end{frame}

\begin{frame}{Estimating Conditional Quantiles}
\begin{itemize}
  \item For a new point $x$, let $L_b(x)$ be the leaf of tree $b$ that $x$ falls into.
  \pause
  \item Define a weight for each training example $i$:
  \begin{equation}
    w_i(x) = \frac{1}{B}\sum_{b=1}^{B}
    \frac{\mathbf{1}\{x_i \in L_b(x)\}}{\bigl|L_b(x)\bigr|},
    \label{eq:qrf-weight}
  \end{equation}
  where $|L_b(x)|$ is the number of training points in that leaf.
  \pause
  \item The weighted empirical CDF is
  \[
    \hat{F}(y\mid x) = \sum_{i=1}^{n} w_i(x)\,\mathbf{1}\{y_i \le y\}.
  \]
  \pause
  \item The $\tau$-quantile prediction is $\hat{Q}_\tau(x) = \inf\{y : \hat{F}(y\mid x) \ge \tau\}$.
\end{itemize}
\end{frame}

\begin{frame}{Worked Numerical Example: Prediction Interval}
\small
\begin{itemize}
  \item A new point $x$ falls into leaves containing these pooled $y$-values (sorted):
  \[
    \{12,\,14,\,15,\,15,\,16,\,17,\,17,\,18,\,19,\,22\}.
  \]
  \pause
  \item Median ($\tau=0.5$) prediction: $\hat{Q}_{0.5}(x) = 16.5$.
  \pause
  \item A $80\%$ prediction interval uses $\tau=0.1$ and $\tau=0.9$:
  \[
    \hat{Q}_{0.1}(x) = 14,\qquad \hat{Q}_{0.9}(x) = 19.
  \]
  \pause
  \item Report: ``predicted value $16.5$, $80\%$ interval $[14, 19]$'' -- far more informative than a single number.
\end{itemize}
\end{frame}

\begin{frame}{Choosing a Forest Variant}
\begin{itemize}
  \item \textbf{Need labeled classification/regression?} Use a standard random forest (Lecture 6).
  \pause
  \item \textbf{Need to flag rare, unusual points with no labels?} Use Isolation Forest.
  \pause
  \item \textbf{Need faster training and lower variance, can tolerate a bit more bias?} Use Extra-Trees.
  \pause
  \item \textbf{Need a prediction interval, not just a point estimate?} Use Quantile Regression Forests.
  \pause
  \item These are not mutually exclusive -- the splitting rule (random vs optimal threshold) and the output rule (vote/mean vs path length vs full distribution) can be mixed.
\end{itemize}
\end{frame}

\begin{frame}{Key Hyperparameters}
\begin{itemize}
  \item \textbf{Isolation Forest:} \texttt{n\_estimators} $B$, subsample size $\psi$ per tree (small, e.g.\ 256), height limit $\approx \log_2 \psi$.
  \pause
  \item \textbf{Extra-Trees:} \texttt{n\_estimators}, \texttt{max\_features} $m$ (as in RF), bootstrap on/off.
  \pause
  \item \textbf{Quantile Regression Forests:} same growth hyperparameters as a regression forest; quantile level $\tau$ is chosen at prediction time, not training time.
\end{itemize}
\end{frame}

\begin{frame}{Summary}
\begin{itemize}
  \item Isolation Forest scores anomalies by how quickly random partitioning isolates a point.
  \item The anomaly score $s(x,\psi) = 2^{-\mathbb{E}[h(x)]/c(\psi)}$ maps short paths to scores near 1.
  \item Extra-Trees randomizes the split threshold itself, trading a little bias for less correlation and faster training.
  \item Quantile Regression Forests keep each leaf's full $y$-distribution to produce prediction intervals, not just a mean.
  \item All three reuse the bagging/feature-randomization machinery from random forests for a different purpose.
\end{itemize}
\end{frame}

\begin{frame}{Lab Connection}
\begin{itemize}
  \item Implement random partitioning and path-length tracking for an isolation tree.
  \pause
  \item Compute $c(\psi)$ and the anomaly score $s(x,\psi)$ for a toy dataset with injected outliers.
  \pause
  \item Modify your Lab 6 split search to draw one random threshold per feature (Extra-Trees) and compare accuracy/variance to the standard forest.
  \pause
  \item Extract per-leaf $y$-values from a regression forest and report a prediction interval instead of a single value.
\end{itemize}
\end{frame}
